\chapter{Fundamentação Teórica}

% TODO: parágrafo de abertura indicando como as seções a seguir se
% articulam e qual o caminho conceitual percorrido.

\section{Transformação digital no setor de food service}
\label{sec:td-foodservice}

% TODO: conceito de transformação digital; impacto específico no setor
% de bares e restaurantes; aceleração causada pela pandemia de COVID-19.
% Apoiar com autores de TD em PMEs e relatórios setoriais (ABRASEL).

\section{Cardápios digitais e QR Code}
\label{sec:cardapios-qr}

% TODO: definição e tipologia (PDF estático, web app, app nativo);
% histórico breve da tecnologia QR Code aplicada ao varejo; arquiteturas
% comuns (link -> web app -> backend de cardápio).

\section{Autoatendimento e experiência do usuário em mobile}
\label{sec:autoatendimento-ux}

% TODO: conceito de self-ordering; heurísticas de usabilidade aplicadas
% a cardápios mobile (Nielsen); acessibilidade (WCAG); fricções típicas
% (QR não escaneia, conexão fraca, navegação confusa).

\section{LGPD aplicada a restaurantes}
\label{sec:lgpd}

% TODO: panorama da LGPD; dados pessoais coletados via cardápio digital
% (e-mail, telefone, geolocalização, comportamento de navegação); base
% legal aplicável; obrigações do controlador.

\section{Trabalhos relacionados}
\label{sec:trabalhos-relacionados}

% TODO: 4--8 trabalhos (TCCs, dissertações, artigos) que tratem de
% digitalização em restaurantes, QR menus ou adoção de tecnologia em PMEs.
% Encerrar com quadro comparativo (foco, método, recorte, lacuna).
